\documentclass[11pt]{article}

\usepackage[margin=1in]{geometry}
\usepackage[T1]{fontenc}
\usepackage[utf8]{inputenc}
\usepackage{lmodern}
\usepackage{booktabs}
\usepackage{microtype}
\usepackage[hidelinks]{hyperref}
\hypersetup{
  pdftitle={Intensity-Clipping Robustness for Compact Lesion Segmentation},
  pdfauthor={Internal manuscript draft}
}

\title{Intensity-Clipping Robustness for Compact Lesion Segmentation}
\author{Internal manuscript draft}
\date{Draft compiled from supplied notes}

\begin{document}
\maketitle

\begin{abstract}
Small-lesion segmentation can be sensitive to intensity variation at inference time. This compact manuscript draft summarizes a preliminary internal validation study in which test-time percentile clipping was added before an existing 3D U-Net inference pipeline. On one internal validation split of 18 volumes, the clipped pipeline increased small-lesion Dice from 0.421 to 0.469 and mean Dice from 0.742 to 0.758, while HD95 changed from 18.4 mm to 17.9 mm. These observations suggest that intensity clipping may improve segmentation stability on this split, but they do not establish broad robustness: a motion-corrupted scan remained a failure case, and no external validation or statistical testing is available in the supplied evidence.
\end{abstract}

\section{Introduction}

Compact lesion segmentation pipelines are often evaluated under constrained validation settings before broader clinical or multi-site testing is available. In this setting, a small change to test-time preprocessing can be useful if it improves lesion delineation without requiring model retraining. The present draft evaluates one such change: percentile-based intensity clipping applied before an existing 3D U-Net inference pipeline.

The evidence is intentionally narrow. The available notes describe one internal validation split with 18 volumes, three aggregate metrics, and one persistent failure case. The manuscript therefore treats the result as preliminary evidence for a scoped preprocessing effect, not as a general robustness claim.

\section{Methods}

\subsection{Validation setting}

The study used one internal validation split containing 18 volumes. The source notes do not specify the dataset name, acquisition protocol, lesion definition, scanner distribution, or train/validation partitioning beyond this split size. These details should be supplied before external submission or archival release.

\subsection{Segmentation pipeline}

The baseline system was an existing 3D U-Net inference pipeline for lesion segmentation. The evaluated modification added test-time percentile clipping before inference. Because the notes do not define the exact lower and upper percentile thresholds, the present source package records the clipping operation as a pending method detail rather than inventing a parameter value.

\subsection{Evaluation metrics}

The supplied evidence reports small-lesion Dice, mean Dice, and the 95th-percentile Hausdorff distance (HD95, in millimeters). No confidence intervals, statistical tests, per-case distributions, or subgroup analyses were provided. Results are therefore presented as descriptive aggregate changes on the internal split.

\section{Results}

Table~\ref{tab:metrics} summarizes the reported aggregate metrics. Percentile clipping was associated with higher Dice values and a small reduction in HD95 on the internal validation split.

\begin{table}[htbp]
\centering
\caption{Reported internal validation metrics for the baseline pipeline and the test-time percentile-clipping variant.}
\label{tab:metrics}
\begin{tabular}{lccc}
\toprule
Metric & Baseline & Percentile clipping & Change \\
\midrule
Small-lesion Dice & 0.421 & 0.469 & +0.048 \\
Mean Dice & 0.742 & 0.758 & +0.016 \\
HD95 (mm) & 18.4 & 17.9 & -0.5 \\
\bottomrule
\end{tabular}
\end{table}

The largest reported change was for small-lesion Dice, which increased by 0.048 absolute points. Mean Dice increased by 0.016, and HD95 decreased by 0.5 mm. These values are consistent with a modest improvement in this internal split. The supplied notes also identify a motion-corrupted scan that remained a failure case after clipping, indicating that the preprocessing change did not resolve all image-quality-driven errors.

\section{Limitations}

The study is preliminary. It uses one internal validation split with 18 volumes and does not include external validation, repeated splits, statistical uncertainty, ablation of clipping thresholds, or per-case error analysis. The exact percentile bounds are not specified in the source notes. The persistent failure on a motion-corrupted scan further limits any broad robustness interpretation. The current evidence supports only a scoped statement that test-time percentile clipping may improve aggregate lesion-segmentation metrics on this internal split.

\section{Discussion}

The observed pattern suggests that percentile clipping is a plausible low-cost inference-time adjustment for this pipeline. The improvement is most visible in small-lesion Dice, which is the metric most likely to reflect sensitivity to small target structures. At the same time, the evidence does not show that clipping generalizes across datasets, scanners, lesion types, or severe artifacts. The motion-corrupted failure case should remain visible in any downstream manuscript version, because it defines the current boundary of the result.

Before this draft is treated as a submission-ready manuscript, the next production pass should add the exact clipping thresholds, dataset and split details, citation-supported background, per-case or uncertainty summaries, and venue-specific formatting. The scientific claim should remain bounded unless additional experiments support a broader conclusion.

\section*{Citation Status}

No bibliography is included in this compact source package. Citation support is pending and must be supplied explicitly before submission or public release; no external citations were invented for this draft.

\end{document}
